\documentclass[12pt, a4paper]{article}
\usepackage[utf8]{inputenc}
\usepackage[T1]{fontenc}
\usepackage[spanish]{babel}
\usepackage{amsmath, amssymb, amsthm}
\usepackage{geometry}
%\rfoot{P\'agina \thepage \hspace{1pt} de 2}
\everymath{\displaystyle}

% Configuración de márgenes para el formato de examen
\geometry{
 a4paper,
 total={170mm,257mm},
 left=20mm,
 top=20mm,
 }

% Título y detalles del examen


\title{\textbf{Examen Parcial de Probabilidad}}
\author{Universidad (Nombre de la Institución) \\ \small Curso: Probabilidad o Probabilidad y Estadística \\ \small Semestre: Tercero}
%\date{Tiempo Máximo: 110 minutos}

\title{Examen Parcial de Probabilidad y Estad\'istica}

\author{Prof. Mtro. Luis Jos\'{e} Yudico Anaya}
\date{06-10-2025}


\begin{document}

\maketitle
\thispagestyle{empty} % No número de página en la primera hoja



\noindent\rule{\linewidth}{0.5pt}

\section*{Instrucciones}
\begin{enumerate}
    \item El examen consta de \textbf{5 preguntas} (1 de concepto y 4 problemas). Resuelva cada uno de ellos mostrando \textbf{todo el procedimiento y justificación teórica} a mano.
    \item La respuesta final para cada inciso debe estar claramente recuadrada.
\end{enumerate}

\noindent\rule{\linewidth}{0.5pt}

\section*{Pregunta Conceptual}
\begin{enumerate}
    \item[\textbf{C-1.}] \textbf{Espacio de Probabilidad}
    \\
    \textbf{Defina} formalmente un \textbf{Espacio de Probabilidad}. Mencione y describa brevemente cada uno de sus tres componentes $(\Omega, \mathcal{F}, \mathbb{P})$, prestando especial atención a las propiedades que debe cumplir el $\sigma$-álgebra $\mathcal{F}$.
\end{enumerate}

\noindent\rule{\linewidth}{0.5pt}

\section*{Problemas}

\begin{enumerate}

    \item[\textbf{1.}] \textbf{Métodos de Conteo en Diseño Experimental de Nanomateriales}
    \\
    Un equipo de investigación en nanotecnología debe diseñar un experimento de síntesis de nanocristales. Tienen a su disposición \textbf{8 catalizadores distintos} ($C_1$ a $C_8$) y \textbf{5 tipos de nanopartículas base} ($N_1$ a $N_5$).
    \\
    Determine la cantidad de formas distintas en que pueden configurarse las siguientes etapas de diseño:
    \begin{enumerate}
        \item[\textbf{i)}] \textbf{Selección de Parámetros de Control (Principio Multiplicativo):} El sistema de control del reactor utiliza un código de 6 dígitos para establecer la temperatura y presión. Si cada dígito puede ser cualquier valor del 1 al 9 y la repetición está permitida, ¿cuántos códigos de control son posibles?
        \item[\textbf{ii)}] \textbf{Orden de Aplicación de Catalizadores (Permutación sin Repetición):} Si se deben probar 3 catalizadores diferentes de los 8 disponibles, y el orden en que se aplican al reactor es crucial para el resultado, ¿cuántas secuencias de aplicación únicas de 3 catalizadores son posibles?
        \item[\textbf{iii)}] \textbf{Mezcla de Nanopartículas Base (Combinación con Repetición):} Para un lote de prueba, se necesita seleccionar una muestra de 4 nanopartículas base. El orden de selección no importa, y se pueden seleccionar varias muestras del mismo tipo. ¿Cuántas combinaciones diferentes de las 4 nanopartículas base son posibles a partir de los 5 tipos disponibles?
    \end{enumerate}

    \vspace{0.8cm}

    \item[\textbf{2.}] \textbf{Axiomas y Corolarios de la Probabilidad}
    \\
    Sean $A$ y $B$ dos eventos en un espacio de probabilidad tales que $\mathbb{P}(A) = 0.5$, $\mathbb{P}(B) = 0.4$, y $\mathbb{P}(A \cap B) = 0.2$.
    \\
    Calcule las siguientes probabilidades, justificando la fórmula o corolario utilizado en cada paso:
    \\
    a) $\mathbb{P}(A \cup B)$
    \\
    b) $\mathbb{P}(A^c)$
    \\
    c) $\mathbb{P}((A \cup B)^c)$
    \\
    d) La probabilidad de que \textbf{solamente} ocurra el evento $B$.

    \vspace{0.8cm}

    \item[\textbf{3.}] \textbf{Probabilidad Condicional e Independencia}
    \\
    Se lanza un dado de seis caras (numeradas del 1 al 6). Considere los siguientes eventos:
    \begin{itemize}
        \item $A$: Se obtiene un número par.
        \item $B$: Se obtiene un número mayor que 3.
        \item $C$: Se obtiene un número primo.
    \end{itemize}
    a) Calcule la probabilidad condicional $\mathbb{P}(B \mid A)$.
    \\
    b) Calcule la probabilidad condicional $\mathbb{P}(A \mid C)$.
    \\
    c) Determine si los eventos $A$ y $B$ son \textbf{independientes}. Justifique su respuesta utilizando la definición formal de independencia.
    \\
    d) Determine si los eventos $A$ y $C$ son \textbf{independientes}. Justifique su respuesta utilizando la definición formal de independencia.

    \vspace{0.8cm}

    \item[\textbf{4.}] \textbf{Teorema de Bayes en Detección de Defectos Nanométricos}
    \\
    Una línea de producción fabrica chips de memoria basados en nanotubos de carbono, utilizando tres lotes de materia prima: \textbf{Lote A} ($\mathbb{P}(A)=30\%$), \textbf{Lote B} ($\mathbb{P}(B)=45\%$) y \textbf{Lote C} ($\mathbb{P}(C)=25\%$). La probabilidad de que un chip sea defectuoso ($D$) depende del lote: $\mathbb{P}(D \mid A) = 1\%$, $\mathbb{P}(D \mid B) = 2\%$, $\mathbb{P}(D \mid C) = 3\%$.
    \\
    \begin{enumerate}
        \item[\textbf{a)}] \textbf{Ley de Probabilidad Total}: Si se selecciona un chip al azar, ¿cuál es la probabilidad de que sea defectuoso ($D$)?
        \item[\textbf{b)}] \textbf{Teorema de Bayes}: Si se selecciona un chip y se detecta que es defectuoso ($D$), ¿cuál es la probabilidad de que provenga del Lote C? Es decir, calcule $\mathbb{P}(C \mid D)$.
        \item[\textbf{c)}] \textbf{Concepto}: ¿Cuál es el término que describe las probabilidades $\mathbb{P}(A)$, $\mathbb{P}(B)$, $\mathbb{P}(C)$ \textbf{antes} de la inspección de defectos? ¿Cuál es el término para $\mathbb{P}(C \mid D)$?
    \end{enumerate}

\end{enumerate}

\end{document}